\documentclass[11pt,letterpaper]{article}

% ── Packages ──────────────────────────────────────────────────────────────────
\usepackage[margin=1in]{geometry}
\usepackage[T1]{fontenc}
\usepackage[utf8]{inputenc}
\usepackage{lmodern}
\usepackage{microtype}
\usepackage{parskip}
\usepackage{titlesec}
\usepackage{fancyhdr}
\usepackage{booktabs}
\usepackage{array}
\usepackage{tabularx}
\usepackage{adjustbox}
\usepackage{pdflscape}
\usepackage{enumitem}
\usepackage{textcomp}          % \textmu
\usepackage{xcolor}
\usepackage{hyperref}
\usepackage{caption}

% ── Hyperref setup ────────────────────────────────────────────────────────────
\hypersetup{
    colorlinks = true,
    linkcolor  = blue!60!black,
    urlcolor   = blue!60!black,
    citecolor  = blue!60!black,
    pdftitle   = {Senior Project Lab Notes},
    pdfauthor  = {},
}

% ── Section formatting ────────────────────────────────────────────────────────
\titleformat{\section}{\large\bfseries}{}{0em}{}[\titlerule]
\titleformat{\subsection}{\normalsize\bfseries}{}{0em}{}
\titleformat{\subsubsection}{\normalsize\itshape}{}{0em}{}

% ── Header / Footer ───────────────────────────────────────────────────────────
\pagestyle{fancy}
\fancyhf{}
\rhead{Senior Project --- Lab Notes}
\lhead{\leftmark}
\rfoot{Page \thepage}
\renewcommand{\headrulewidth}{0.4pt}

% ── Convenience macros ────────────────────────────────────────────────────────
\newcommand{\us}{\textmu{}S}   % microseconds symbol
\newcommand{\code}[1]{\texttt{#1}}
\newcommand{\na}{\textit{N/A}}
\newcommand{\dash}{---}

% ── Column types ──────────────────────────────────────────────────────────────
\newcolumntype{L}[1]{>{\raggedright\arraybackslash}p{#1}}
\newcolumntype{R}[1]{>{\raggedleft\arraybackslash}p{#1}}
\newcolumntype{C}[1]{>{\centering\arraybackslash}p{#1}}

% ══════════════════════════════════════════════════════════════════════════════
\begin{document}

\begin{center}
    {\LARGE\bfseries Senior Project --- Lab Notes}\\[0.4em]
    {\normalsize Last revised: March 1, 2026}
\end{center}

\tableofcontents
\newpage

% ══════════════════════════════════════════════════════════════════════════════
\section{System Environment}
% ══════════════════════════════════════════════════════════════════════════════

\subsection{Broker / Verifier}

\begin{tabular}{@{}L{5cm}L{9cm}@{}}
\toprule
\textbf{Property} & \textbf{Value} \\
\midrule
OS Name            & Microsoft Windows 11 Home \\
OS Version         & 10.0.26100 N/A Build 26100 \\
OS Configuration   & Standalone Workstation \\
OS Build Type      & Multiprocessor Free \\
System Manufacturer & HP \\
System Model       & OMEN by HP Laptop PC \\
System Type        & x64-based PC \\
Processor          & AMD64 Family 23 Model 96 Stepping 1, AuthenticAMD \textasciitilde{}2900\,MHz \\
BIOS Version       & AMI F.13, 3/4/2021 \\
Total Physical Memory & 15,731\,MB \\
Available Physical Memory & 6,761\,MB \\
Virtual Memory Max & 19,827\,MB \\
Network Card       & Intel\textsuperscript{\textregistered} Wi-Fi 6 AX200 160MHz \\
\bottomrule
\end{tabular}

\subsection{Thermometer Device (ESP32)}

\begin{tabular}{@{}L{5cm}L{9cm}@{}}
\toprule
\textbf{Property} & \textbf{Value} \\
\midrule
Chip Model    & ESP32-D0WD-V3 \\
Chip Revision & 301 \\
CPU Frequency & 240\,MHz \\
Flash Size    & 4\,MB \\
Flash Speed   & 80\,MHz \\
Free Heap     & 287,160 bytes \\
SDK Version   & v5.5.1-931-g9bb7aa84fe \\
\bottomrule
\end{tabular}

\subsection{Smart Lights Device (Pi 3B+)}

\subsubsection{CPU}

\begin{tabular}{@{}L{5cm}L{9cm}@{}}
\toprule
\textbf{Property} & \textbf{Value} \\
\midrule
Architecture    & aarch64 \\
CPU op-modes    & 32-bit, 64-bit \\
Byte Order      & Little Endian \\
CPU(s)          & 4 (on-line: 0--3) \\
Vendor ID       & ARM \\
Model name      & Cortex-A53 \\
Threads per core & 1 \\
Cores per cluster & 4 \\
CPU scaling MHz & 43\% \\
CPU max MHz     & 1400.0 \\
CPU min MHz     & 600.0 \\
BogoMIPS        & 38.40 \\
Flags           & fp asimd evtstrm crc32 cpuid \\
L1d cache       & 128\,KiB (4 instances) \\
L1i cache       & 128\,KiB (4 instances) \\
L2 cache        & 512\,KiB (1 instance) \\
NUMA node(s)    & 1 (node0: CPU 0--3) \\
\bottomrule
\end{tabular}

\subsubsection{Vulnerability Mitigations}

All vulnerabilities listed as \textit{Not affected} with the following exception:
\begin{itemize}[noitemsep]
    \item Spectre v1: Mitigation; \code{\_\_user} pointer sanitization
\end{itemize}

\subsubsection{Kernel, Memory, and Storage}

\begin{tabular}{@{}L{5cm}L{9cm}@{}}
\toprule
\textbf{Property} & \textbf{Value} \\
\midrule
Kernel Version   & Linux pi3 6.12.47+rpt-rpi-v8 \#1 SMP PREEMPT Debian 1:6.12.47-1+rpt1 (2025-09-16) aarch64 \\
Total Memory     & 906\,MiB \\
Used Memory      & 149\,MiB \\
Free Memory      & 662\,MiB \\
Swap Total       & 905\,MiB \\
Swap Used        & 0\,MiB \\
Storage Name     & SC16G \\
Storage Speed    & \code{0x0235844300000000} \\
Available Storage & 11\,G of 15\,G (\code{/dev/mmcblk0p2}) \\
Pi 3B+ Revision  & a020d4 \\
\bottomrule
\end{tabular}

\subsubsection{Baseline Task Snapshot}

\begin{verbatim}
top - 22:25:01 up 41 min, 1 user, load average: 0.00, 0.00, 0.00
Tasks: 143 total, 1 running, 142 sleeping, 0 stopped, 0 zombie
%Cpu(s): 0.1 us, 0.2 sy, 0.0 ni, 99.8 id, 0.0 wa, 0.0 hi, 0.0 si, 0.0 st
MiB Mem: 906.0 total, 659.5 free, 150.5 used, 152.3 buff/cache
MiB Swap: 906.0 total, 906.0 free, 0.0 used, 755.6 avail
\end{verbatim}

\subsection{Camera Device (Pi 5)}

\subsubsection{CPU}

\begin{tabular}{@{}L{5cm}L{9cm}@{}}
\toprule
\textbf{Property} & \textbf{Value} \\
\midrule
Architecture    & aarch64 \\
CPU op-modes    & 32-bit, 64-bit \\
Byte Order      & Little Endian \\
CPU(s)          & 4 (on-line: 0--3) \\
Vendor ID       & ARM \\
Model name      & Cortex-A76 \\
Threads per core & 1 \\
Cores per cluster & 4 \\
CPU scaling MHz & 62\% \\
CPU max MHz     & 2400.0 \\
CPU min MHz     & 1500.0 \\
BogoMIPS        & 108.00 \\
Flags           & fp asimd evtstrm aes pmull sha1 sha2 crc32 atomics fphp asimdhp \ldots \\
L1d cache       & 256\,KiB (4 instances) \\
L1i cache       & 256\,KiB (4 instances) \\
L2 cache        & 2\,MiB (4 instances) \\
L3 cache        & 2\,MiB (1 instance) \\
NUMA node(s)    & 8 (nodes 0--7, each: CPU 0--3) \\
\bottomrule
\end{tabular}

\subsubsection{Vulnerability Mitigations}

Most vulnerabilities listed as \textit{Not affected}. Active mitigations:
\begin{itemize}[noitemsep]
    \item Spec store bypass: Mitigation; Speculative Store Bypass disabled via prctl
    \item Spectre v1: Mitigation; \code{\_\_user} pointer sanitization
    \item Spectre v2: Mitigation; CSV2, BHB
\end{itemize}

\subsubsection{Kernel, Memory, and Storage}

\begin{tabular}{@{}L{5cm}L{9cm}@{}}
\toprule
\textbf{Property} & \textbf{Value} \\
\midrule
Kernel Version   & Linux pi5 6.12.47+rpt-rpi-2712 \#1 SMP PREEMPT Debian 1:6.12.47-1+rpt1 (2025-09-16) aarch64 \\
Total Memory     & 8,063\,MiB \\
Used Memory      & 253\,MiB \\
Free Memory      & 7,635\,MiB \\
Swap Total       & 2,047\,MiB \\
Swap Used        & 0\,MiB \\
Storage Name     & SDABC \\
Storage Speed    & \code{02b5800300000000} \\
Available Storage & 8.8\,G of 14\,G (\code{/dev/mmcblk0p2}) \\
Pi 5 Revision    & d04171 \\
\bottomrule
\end{tabular}

\subsubsection{Baseline Task Snapshot}

\begin{verbatim}
top - 22:31:15 up 47 min, 1 user, load average: 0.00, 0.00, 0.00
Tasks: 162 total, 1 running, 161 sleeping, 0 stopped, 0 zombie
%Cpu(s): 0.0 us, 0.0 sy, 0.0 ni, 100.0 id, 0.0 wa, 0.0 hi, 0.0 si, 0.0 st
MiB Mem: 8063.0 total, 7630.9 free, 257.2 used, 256.9 buff/cache
MiB Swap: 2048.0 total, 2048.0 free, 0.0 used, 7805.8 avail
\end{verbatim}

% ══════════════════════════════════════════════════════════════════════════════
\section{MQTT Configuration}
% ══════════════════════════════════════════════════════════════════════════════

\begin{tabular}{@{}L{5cm}L{9cm}@{}}
\toprule
\textbf{Device} & \textbf{Library} \\
\midrule
Thermometer (ESP32) & \code{PubSubClient.h} \\
Smart Lights (Pi 3) & \code{paho.mqtt.client} \\
Camera (Pi 5)       & \code{paho.mqtt.client} \\
Broker / Verifier   & \code{paho.mqtt.client} \\
\bottomrule
\end{tabular}

% ══════════════════════════════════════════════════════════════════════════════
\section{Baseline Test Environment}
% ══════════════════════════════════════════════════════════════════════════════

With no configured background tasks running, baseline tests were conducted to
measure latency on each IoT device.  This included:
\begin{itemize}[noitemsep]
    \item \textbf{Sign latency} --- the additional time required to perform an
          Ed25519 signature calculation on-device.
    \item \textbf{Verify latency} --- the additional time required by the broker
          to perform Ed25519 signature verification (device-level signing
          operations only).
\end{itemize}

% ══════════════════════════════════════════════════════════════════════════════
\section{Device-Level Signing Results}
% ══════════════════════════════════════════════════════════════════════════════

\subsection{ESP32 Thermometer (Broker Verification)}

\textbf{Source:} \code{Broker/ESP32/device\_level\_signing/results/benchmark\_results.csv}\\
\textbf{Total entries:} 10,000 \quad \textbf{Success rate:} 100.00\%

\begin{tabular}{@{}L{6cm}R{4cm}@{}}
\toprule
\textbf{Metric} & \textbf{Value} \\
\midrule
Avg Sign Time (\us) & 46,319.87 \\
Max Sign Time (\us) & 46,393.00 \\
Avg Verify Time (\us) & 202.94 \\
Max Verify Time (\us) & 981.80 \\
\bottomrule
\end{tabular}

\subsection{Pi3 Smart Lights (Broker Verification)}

\begin{adjustbox}{max width=\linewidth}
\begin{tabular}{@{}L{3cm}L{8.5cm}R{1.6cm}R{1.8cm}R{1.8cm}R{1.8cm}R{1.8cm}R{1.8cm}@{}}
\toprule
\textbf{Scenario} & \textbf{Source File} & \textbf{Entries} & \textbf{Success} &
\textbf{Avg Sign} & \textbf{Max Sign} & \textbf{Avg Verify} & \textbf{Max Verify} \\
 & & & & \multicolumn{4}{c}{(\us)} \\
\midrule
Early Dev Run &
  \code{Broker/Pi3/device\_level\_signing/} \code{benchmark\_pi\_results\_1.csv} &
  5,000 & 100.00\% & 99,192.21 & 194,199.00 & 36,497.29 & 56,841.50 \\
No Stress &
  \code{Broker/Pi3/device\_level\_signing/} \code{results/benchmark\_pi\_results.csv} &
  5,000 & 100.00\% & 487.07 & 910.00 & 239.83 & 817.80 \\
Stress 50 &
  \code{Broker/Pi3/device\_level\_signing/} \code{results/benchmark\_pi\_results\_(stress-50).csv} &
  5,000 & 100.00\% & 294.40 & 5,265.00 & 212.41 & 10,549.70 \\
Stress 99 &
  \code{Broker/Pi3/device\_level\_signing/} \code{results/benchmark\_pi\_results\_(stress-99).csv} &
  5,000 & 100.00\% & 339.17 & 4,800.00 & 216.40 & 855.70 \\
\bottomrule
\end{tabular}
\end{adjustbox}

\textit{Note: The Early Dev Run file is located outside the \code{results/} subdirectory and predates
the optimised implementation; sign times are approximately 200$\times$ higher than production runs.}

\subsection{Pi5 Camera (Device-Level Signing Benchmark Summaries)}

\begin{adjustbox}{max width=\linewidth}
\begin{tabular}{@{}L{2cm}L{7.5cm}L{3.5cm}R{1.5cm}R{1.5cm}R{1.5cm}R{1.5cm}R{1.8cm}R{1.8cm}@{}}
\toprule
\textbf{Scenario} & \textbf{Source File} & \textbf{Timestamp} &
\textbf{Chunks} & \textbf{Success} &
\textbf{Avg Sign} & \textbf{Max Sign} & \textbf{Avg Verify} & \textbf{Max Verify} \\
 & & & & & \multicolumn{4}{c}{(\us)} \\
\midrule
No Stress &
  \code{Broker/Pi5/device\_level\_signing/} \code{results/no~stress/benchmark\_summary.csv} &
  2026-01-26 16:30:34 & 4,238 & \na & 165.53 & 3,465 & 309.94 & 913.70 \\
Stress 25 &
  \code{Broker/Pi5/device\_level\_signing/} \code{results/stress25/benchmark\_summary\_1.csv} &
  2026-02-07 16:09:32 & 4,230 & 100.00\% & 139.97 & 5,161 & 293.77 & 5,067.70 \\
Stress 50 &
  \code{Broker/Pi5/device\_level\_signing/} \code{results/stress50/benchmark\_summary1.csv} &
  2026-02-07 16:01:24 & 4,377 & \na & 149.95 & 7,629 & 282.89 & 10,570.00 \\
Stress 75 &
  \code{Broker/Pi5/device\_level\_signing/} \code{results/stress75/benchmark\_summary\_1.csv} &
  2026-02-07 16:11:23 & 4,230 & 100.00\% & 158.46 & 5,565 & 292.48 & 1,257.90 \\
Stress 99 &
  \code{Broker/Pi5/device\_level\_signing/} \code{results/stress99/benchmark\_summary\_1.csv} &
  2026-02-07 16:13:00 & 4,228 & 100.00\% & 173.04 & 7,376 & 286.30 & 1,385.70 \\
\bottomrule
\end{tabular}
\end{adjustbox}

% ══════════════════════════════════════════════════════════════════════════════
\section{IPFS Results}
% ══════════════════════════════════════════════════════════════════════════════

\subsection{Pi5 IPFS Hash + Verify}

\begin{adjustbox}{max width=\linewidth}
\begin{tabular}{@{}R{0.6cm}L{7.5cm}L{3.5cm}R{1.5cm}R{1.5cm}R{1.8cm}R{1.8cm}R{1.8cm}R{1.8cm}@{}}
\toprule
\textbf{Run} & \textbf{Source File} & \textbf{Timestamp} &
\textbf{Chunks} & \textbf{Success} &
\textbf{Avg Hash} & \textbf{Max Hash} & \textbf{Avg Verify} & \textbf{Max Verify} \\
 & & & & & \multicolumn{4}{c}{(\us)} \\
\midrule
1 & \code{Broker/Pi5/IPFS/results/benchmark\_summary\_1.csv} & 2026-02-15 22:01:42 &
  2,481 & 100.00\% & 19,855.99 & 32,443.00 & 36,607.57 & 49,253.40 \\
2 & \code{Broker/Pi5/IPFS/results/benchmark\_summary\_2.csv} & 2026-02-15 22:09:16 &
  2,692 & 100.00\% & 19,913.55 & 29,267.00 & 36,743.69 & 47,372.20 \\
\bottomrule
\end{tabular}
\end{adjustbox}

\subsection{ESP32 \texorpdfstring{$\rightarrow$}{->} Broker/Kubo IPFS Pipeline}

Runs 1--2 predate the P50/P95/P99 schema; those cells are marked ``\dash''.

\begin{landscape}
\begin{adjustbox}{max width=\linewidth}
\begin{tabular}{@{}R{0.5cm}
                L{5.5cm}
                L{3.2cm}
                R{1.3cm}R{1.4cm}
                R{1.5cm}R{1.4cm}
                R{1.6cm}R{1.8cm}R{1.5cm}R{1.5cm}
                R{1.6cm}R{1.8cm}R{1.5cm}R{1.5cm}@{}}
\toprule
\textbf{Run} & \textbf{Source File} & \textbf{Timestamp} &
\textbf{Chunks} & \textbf{Success} &
\multicolumn{2}{c}{\textbf{Device (\us)}} &
\multicolumn{4}{c}{\textbf{Hash (\us)}} &
\multicolumn{4}{c}{\textbf{Verify (\us)}} \\
\cmidrule(lr){6-7}\cmidrule(lr){8-11}\cmidrule(lr){12-15}
& & & & & \textbf{Avg} & \textbf{Max} &
\textbf{Avg} & \textbf{Max} & \textbf{P50} & \textbf{P95} &
\textbf{Avg} & \textbf{Max} & \textbf{P50} & \textbf{P95} \\
\midrule
1 & \code{Broker/ESP32/IPFS/results/} \code{benchmark\_summary\_1.csv} & 2026-03-01 19:43:00 &
  52 & 100.00\% & 448.27 & 495 & 62,062.56 & 382,014.00 & \dash & \dash & 55,782.32 & 62,952.70 & \dash & \dash \\

2 & \code{Broker/ESP32/IPFS/results/} \code{benchmark\_summary\_2.csv} & 2026-03-01 19:54:38 &
  5,939 & 100.00\% & 449.05 & 568 & 56,378.79 & 293,549.80 & \dash & \dash & 56,580.27 & 716,918.90 & \dash & \dash \\

3 & \code{Broker/ESP32/IPFS/results/} \code{benchmark\_summary\_3.csv} & 2026-03-01 20:33:49 &
  772 & 76.94\% & 450.93 & 533 & 56,482.97 & 88,893.90 & 55,605.30 & 62,254.20 & 52,269.25 & 87,903.50 & 54,930.50 & 58,613.30 \\

4 & \code{Broker/ESP32/IPFS/results/} \code{benchmark\_summary\_4.csv} & 2026-03-01 20:41:47 &
  119 & 91.60\% & 448.18 & 516 & 216,208.58 & 1,137,796.30 & 60,431.80 & 597,412.50 & 245,431.04 & 1,618,677.20 & 60,155.50 & 618,591.90 \\

5 & \code{Broker/ESP32/IPFS/results/} \code{benchmark\_summary\_5.csv} & 2026-03-01 20:54:20 &
  5,029 & 100.00\% & 448.51 & 529 & 57,599.96 & 383,334.70 & 55,252.50 & 66,373.20 & 57,684.23 & 648,736.60 & 55,292.60 & 66,472.40 \\
\bottomrule
\end{tabular}
\end{adjustbox}
\end{landscape}

\subsection{IPFS Notes}

\begin{itemize}[noitemsep]
    \item Pi5 IPFS runs show stable hash/verify timing around \textasciitilde{}20\,ms /
          \textasciitilde{}36\,ms.
    \item ESP32 pipeline runs \#2 and \#5 provide long-duration datasets (5,939 and 5,029
          chunks respectively); device-side prep remained near \textasciitilde{}449\,\us{}
          across all runs while broker-side Kubo hash/verify dominated latency.
    \item Run \#3 (76.94\% success rate) reflects packet loss at the MQTT ingress layer,
          not a signing or hashing failure.
    \item Run \#4 (119 chunks, 91.60\% success) shows extreme broker-side latency spikes
          (Max Hash: 1,137,796\,\us{}; Max Verify: 1,618,677\,\us{}); P50 values
          (\textasciitilde{}60\,ms) indicate these are outliers, not representative of
          steady-state performance.
    \item Runs \#3--\#5 include P50/P95/P99 percentile columns in the source CSVs;
          runs \#1--\#2 predate that schema and report avg/max only.
\end{itemize}

% ══════════════════════════════════════════════════════════════════════════════
\section{System Environment --- Confirmation Record}
% ══════════════════════════════════════════════════════════════════════════════

The system environment block in Section~1 is the canonical environment record for:

\begin{itemize}[noitemsep]
    \item Broker/Verifier (Windows 11 laptop)
    \item Thermometer Device (ESP32)
    \item Smart Lights Device (Pi 3)
    \item Camera Device (Pi 5)
\end{itemize}

Additional ESP32 flash/run confirmation captured during the latest IPFS Arduino deployment:

\begin{itemize}[noitemsep]
    \item Board: \code{esp32:esp32:esp32}
    \item Chip: ESP32-D0WD-V3 (revision v3.1)
    \item Upload port: COM5
    \item Flash/upload completed successfully via esptool.
\end{itemize}

% ══════════════════════════════════════════════════════════════════════════════
\section{Special Consideration Notes}
% ══════════════════════════════════════════════════════════════════════════════

There is currently no mature, production-ready IPFS node implementation for ESP32
that can run the full Kubo/libp2p stack (peer discovery, block exchange, datastore,
and networking services) directly on-device.  This is primarily a hardware and
software ecosystem limitation: ESP32 memory, storage, and CPU constraints are not
well matched to full IPFS node requirements.

Because of this, the ESP32 was used as a constrained edge publisher, and a
broker/gateway host performed true IPFS operations using Kubo
(\code{ipfs add -{}-only-hash}) on the received payload.  This preserves real IPFS
CID generation while still allowing device-side timing and stress behaviour to be
measured.

\subsection*{Interpretation Notes for Results}

\begin{itemize}[noitemsep]
    \item ESP32 device timing reflects sensor read + payload construction + MQTT
          publish overhead.
    \item IPFS hash/verify timing reflects host-side Kubo and verifier execution,
          not native ESP32 hashing.
    \item End-to-end integrity remains valid because CIDs are generated and checked
          by a real IPFS implementation.
    \item This architecture should be reported as \textbf{ESP32 + IPFS gateway
          orchestration}, not standalone ESP32-native IPFS.
\end{itemize}

\subsection*{Additional Benchmark Caveats}

\begin{enumerate}

\item \textbf{Workload mismatch across IPFS runs}

  \textit{Concern:} ESP32 IPFS payloads are fixed at 25 bytes, while Pi5 IPFS
  chunks are mostly \textasciitilde{}4\,KB.  Direct latency comparisons are
  therefore not apples-to-apples.

  \textit{Applies to:} Cross-platform IPFS comparison (ESP32 IPFS pipeline
  vs.\ Pi5 IPFS pipeline), not within-run comparisons on the same device.

  \textit{Evidence:} \code{Broker/ESP32/IPFS/results/raw\_packet\_data\_2.csv}
  and \code{Broker/Pi5/IPFS/results/raw\_packet\_data\_2.csv}.

\item \textbf{Outlier sensitivity in ESP32 IPFS hash/verify time}

  \textit{Concern:} ESP32 IPFS run data has extreme max latency spikes
  (run \#2 Max Verify: 716,918.9\,\us{}; run \#4 Max Hash: 1,137,796.3\,\us{},
  Max Verify: 1,618,677.2\,\us{}; run \#5 Max Verify: 648,736.6\,\us{}),
  so averages alone are highly misleading.

  \textit{Applies to:} ESP32 IPFS benchmark interpretation and any summary
  table that uses only avg/max metrics.

  \textit{Recommendation:} Use P50/P95 as the primary performance indicators;
  report avg/max as supplementary range data only.

  \textit{Evidence:} \code{Broker/ESP32/IPFS/results/benchmark\_summary\_2.csv},
  \code{benchmark\_summary\_4.csv}, \code{benchmark\_summary\_5.csv}.

\item \textbf{Double-hash measurement effect}

  \textit{Concern:} In default validation mode (\code{double\_hash}),
  \code{broker\_ipfs.py} computes CID once for publish and a second time for
  verification, which increases host compute cost versus a single-pass pipeline.

  \textit{Applies to:} ESP32 $\rightarrow$ broker/Kubo IPFS pipeline only
  (the host running \code{Broker/ESP32/IPFS/broker\_ipfs.py}), not the ESP32
  microcontroller itself.

  \textit{Note:} This can be switched to \code{single\_hash} mode when
  throughput-focused benchmarking is needed.

  \textit{Evidence:} \code{Broker/ESP32/IPFS/broker\_ipfs.py}.

\item \textbf{MQTT reliability mode (mixed QoS path)}

  \textit{Concern:} The ESP32 sender path uses QoS~0 (possible packet loss
  without retransmit), while broker republish uses QoS~1; the reported success
  rate reflects received/validated packets, not guaranteed delivery of all
  packets generated at device origin.

  \textit{Applies to:} ESP32 $\rightarrow$ broker ingress reliability
  interpretation and end-to-end delivery claims.

  \textit{Evidence:} \code{ESP32/ipfs/esp32\_ipfs/esp32\_ipfs.ino} and
  \code{Broker/ESP32/IPFS/broker\_ipfs.py}.

\end{enumerate}

\end{document}
